% Philippica X (philippic-10) --- English translation
% Translated by Claude (Anthropic) from the Latin (Perseus).
% Style guide: STYLE.md.

\ciceroSection{1} We all owe you the greatest thanks, Pansa, both to feel and to render --- you who, when we did not suppose you would be holding the Senate today, the moment you received the letter of Marcus Brutus, that most excellent citizen, allowed not the least delay before we could enjoy, as soon as could be, the greatest joy and rejoicing. Your act ought to be welcome to all; but even more so the speech you made once the letter had been read out. For you made plain the truth of what I have always felt: that no one who trusts in his own virtue envies the virtue of another.

\ciceroSection{2} And so I, who am bound to Brutus by the most numerous good offices and by the closest intimacy, have the less to say of him. For the part I had taken up for myself, your speech has forestalled. But upon me, senators, the opinion of the man who was asked before me has laid the necessity of saying a little more --- a man from whom I so often dissent that I now fear that our perpetual disagreement may seem to lessen our friendship, which is the last thing that ought to happen.

\ciceroSection{3} For what is this reasoning of yours, Calenus, what is this cast of mind, that never since the Kalends of January have you been of the same opinion as the man who calls on you to speak first, and never has the Senate been so crowded that any single man has followed your motion? Why do you always defend causes unlike yourself? Why, when your own life and fortune invite you to peace, to dignity, do you approve, decree, and hold opinions hostile both to the common peace and to your own dignity?

\ciceroSection{4} For --- to pass over earlier matters --- this at least, which moves my greatest astonishment, I will not pass over in silence. What war have you with the Bruti? Why do you alone assail those whom nearly all of us are bound to revere? The one being besieged you bear without distress; the other you strip, by your motion, of those forces which he himself, by his own toil and peril, raised for the protection of the commonwealth and not of himself, by his own effort, with no one helping. What is this feeling of yours, this way of thinking --- that you should not approve the Bruti but should approve the Antonii; that those whom all men hold most dear, you should hate, and those whom all others hate most bitterly, you should most steadfastly love? Yours are the most ample fortunes, the highest rank of honour, a son --- born, as I both hear and hope, for glory; and I favour him both for the commonwealth's sake, and also for yours.

\ciceroSection{5} I ask, then: would you rather have him like Brutus or like Antony? --- and I leave you free to choose which of the three Antonii you please. ``The gods forbid!'' you will say. Why, then, do you not favour those, do you not praise those, like whom you wish your son to be? For so you will at once take counsel for the commonwealth and set before him examples to imitate. But on this point, Quintus Fufius, I wish to make my complaint to you --- without offence to our friendship --- as a senator who dissents from you. For you said, and indeed from a written text (lest I should think you had slipped for want of a word), that the letter of Brutus seemed ``well and properly written.'' What is this but to praise Brutus's clerk, and not Brutus?

\ciceroSection{6} Experience in public affairs, Calenus, you both ought to have, and can have, in large measure by now. When have you seen it so decreed --- or by what decree of the Senate of this kind, and they are beyond counting, has the Senate ever pronounced a letter ``well written''? This word did not slip from you by accident, as often happens: you brought it written down, pondered, thought out. If anyone could rid you of this habit of carping at most good causes, what that any man might wish for himself would not be left to you? Therefore collect yourself, and at last calm and soften that temper of yours: listen to good men, of whom you have many about you; talk with that wisest of men, your son-in-law, more often than with yourself --- then at last you will hold the name of the highest honour worthily. Or do you count it as nothing --- a thing for which, out of friendship, I am wont to grieve on your behalf --- that it gets carried abroad and reaches the ears of the Roman people, that to the man who spoke his opinion first not one man assented? Which I think will be the case today as well. You are taking the legions away from Brutus. Which legions? Those, surely, which he turned aside from the crime of Gaius Antonius and brought over to the commonwealth by his own authority. So once again you wish him to be seen stripped bare and, alone, banished from the commonwealth.

\ciceroSection{7} But you, senators, if you abandon and betray Marcus Brutus, what citizen will you ever honour, whom will you favour? Unless perhaps you think that those who set the diadem on a man's head are to be preserved, and those who did away with the very name of kingship are to be cast off. And of that divine and immortal glory of Brutus I will say nothing --- the glory which, enshrined in the most grateful memory of all citizens, has not yet been attested by public authority. Such forbearance, good gods! such restraint, such calm and modesty under injury! He, while he was urban praetor, kept away from the city, dispensed no justice --- though he had recovered all justice for the commonwealth; and though he might have been hedged about by the daily concourse of all good men, which used to flock to him in wondrous number, and by the protection of all Italy, he chose rather to be defended in his absence by the judgment of good men than present by their armed hand: he who did not even put on in person the Games of Apollo, mounted in keeping with his own dignity and that of the Roman people, lest he should open any path to the audacity of the most criminal of men.

\ciceroSection{8} And yet what games, what days, were ever more joyful than when, at every single verse, the Roman people pursued the memory of Brutus with the loudest shouting and applause? The body of the liberator was absent, the memory of liberty was present: and in it the image of Brutus seemed to be discerned. But I saw him myself, in those very days of the games, on the estate of that most distinguished young man Gnaeus Lucullus, his kinsman, thinking of nothing but the peace and concord of the citizens. And I saw the same man afterwards at Velia, withdrawing from Italy so that no occasion for civil war might arise on his account. O that sight, mournful not to men only but to the very waves and shores! that the saviour of his country should withdraw from it, while its destroyers remained within it! The fleet of Cassius followed a few days later, so that I was ashamed, senators, to return to that city from which those men were departing. But with what purpose I returned you heard at the outset, and afterwards you have learned by proof. And so Brutus bided his time.

\ciceroSection{9} For as long as he saw you enduring everything, he himself used an incredible patience: but once he perceived that you were roused toward liberty, he prepared defences for your liberty. And what a plague, and how great, he withstood! For if Gaius Antonius had been able to accomplish what he had set his mind upon --- and he would have been able, had not the virtue of Marcus Brutus stood in the way of his crime --- we should have lost Macedonia, Illyricum, and Greece; Greece would have been either a refuge for Antony if he were driven out, or a rampart for the assault of Italy. Whereas now, under the command, the authority, and the forces of Marcus Brutus --- not merely equipped but even adorned --- Greece stretches out her right hand to Italy and pledges Italy her protection. Whoever takes the army away from him robs the commonwealth both of a most glorious recourse and of a most steadfast protection.

\ciceroSection{10} For my part, I desire that Antony hear this as soon as may be, so that he may understand that it is not Decimus Brutus, whom he besieges with his rampart, but he himself who is besieged. He holds three towns in the whole world; he has Gaul most hostile to him; even those in whom he trusted, the Transpadanes, utterly estranged; all Italy is set against him; the foreign nations, from the first shore of Greece all the way to Egypt, are held by the commands and garrisons of the best and bravest citizens. He had one hope, in Gaius Antonius --- who, set in age midway between his two brothers, vied with each of them in vices. He ran off to Macedonia as though he were being thrust out by the Senate, and not, on the contrary, forbidden to set out.

\ciceroSection{11} What a storm, immortal gods, what a blaze, what devastation, what a plague for Greece, had not the incredible and divine virtue of Brutus crushed the attempt and audacity of that raving man! What speed was his, what care, what virtue! --- though not even the speed of Gaius Antonius is to be despised: had not windfall inheritances delayed him on the road, you would say he had flown, not made a journey. When we wish other men to go to public business, we can scarcely thrust them out: this one we thrust out in the very act of holding him back. But what had he to do with Apollonia, what with Dyrrachium, what with Illyricum, what with the army of the commander Publius Vatinius? He was succeeding Hortensius, as he himself said. The bounds of Macedonia were fixed, his terms fixed, his army fixed --- if indeed he had any at all: but with Illyricum, and with the legions of Vatinius, what had Antonius to do?

\ciceroSection{12} ``But neither had Brutus'' --- for so perhaps some unprincipled man might say. All the legions, all the forces that are anywhere, belong to the commonwealth: nor will those legions which abandoned Marcus Antonius be said to have belonged to Antony rather than to the commonwealth. For every right both to an army and to command is forfeited by the man who uses that command and army to assail the commonwealth. And if the commonwealth herself were to judge, or if all right were settled by her decrees, would she adjudge the legions of the Roman people to Antony, or to Brutus? The one had flown down suddenly to the plunder and ruin of the allies, so that wherever he went he laid all waste, plundered, carried off, and used the army of the Roman people against the Roman people itself; the other had set himself this rule: that wherever he came, a kind of light and hope of deliverance should seem to have come. In short, the one sought forces to overthrow the commonwealth, the other to preserve it. Nor indeed did we see this more clearly than the soldiers themselves, of whom so great a prudence in judging was not to be demanded.

\ciceroSection{13} He writes that Antonius is at Apollonia with seven cohorts --- a man who by now either has been captured (which may the gods grant!) or, at the least, modest fellow that he is, does not enter Macedonia, lest he seem to have acted against the decree of the Senate. A levy has been held in Macedonia by the utmost zeal and industry of Quintus Hortensius, whose excellent spirit, worthy of himself and of his ancestors, you have been able to discern from the letter of Brutus. The legion which Lucius Piso, a legate of Antony, was leading gave itself up to my son Cicero. The cavalry that was being led into Syria, divided in two, the one part abandoned in Thessaly the quaestor by whom it was being led and took itself off to Brutus; the other, in Macedonia, the young Gnaeus Domitius --- a man of the highest virtue, gravity, and constancy --- led away from the Syrian legate. And Publius Vatinius, who was both rightly praised by you before and now deserves praise as well, opened the gates of Dyrrachium to Brutus and handed over his army.

\ciceroSection{14} The commonwealth, then, holds Macedonia, holds Illyricum, guards Greece: ours are the legions, ours the light-armed troops, ours the cavalry, and most of all Brutus is ours, and ever ours --- born for the commonwealth both by his own most excellent virtue and by a kind of destiny of his line and name on his father's side and his mother's. Does anyone, then, fear war from this man --- who, before we were forced to take it up, preferred to lie idle in peace rather than to flourish in war? And yet he never lay idle, nor can that word fall upon so great an excellence of virtue. For he was in the longing of the state, on its lips, in the talk of all; yet so far was he from war that, when Italy was ablaze with desire for liberty, he failed the zeal of the citizens rather than lead them into the hazard of arms. And so those very men, if there are any, who blame the slowness of Brutus nonetheless admire his restraint and forbearance.

\ciceroSection{15} But by now I see what they are saying; for they do not do it in secret. They say they are afraid how the veterans will take it that Brutus has an army. As though, indeed, there were any difference between the army of Aulus Hirtius, of Gaius Pansa, of Decimus Brutus, of Gaius Caesar, and this army of Marcus Brutus. For if the four armies I have named are praised because they took up arms for the liberty of the Roman people, what reason is there why this army of Marcus Brutus should not be placed in the same cause? ``But the name of Marcus Brutus is suspect to the veterans.'' More so than that of Decimus? I, for one, do not think it. For although the deed of the Bruti is shared, and the partnership in glory equal, yet those who grieved at that deed were the angrier with Decimus, the more they said that the thing ought least to have been done by him. What, then, are so many armies doing now, except that Brutus may be freed from the siege? And who lead these armies? Those, I suppose, who wish the acts of Gaius Caesar overthrown, who wish the cause of the veterans betrayed.

\ciceroSection{16} If Gaius Caesar himself were alive, would he, I suppose, defend his own acts more keenly than that most valiant man Hirtius defends them? or can anyone be found more friendly to the cause than his son? Yet of these two, the one, not yet recovered from the long course of a most grave illness, brought all the strength he had to the defence of the liberty of those by whose prayers he judged he had been called back from death; the other, stronger in the sturdiness of his virtue than of his years, set out with those very veterans to free Decimus Brutus. So those most certain and most ardent champions of Caesar's acts are waging war for the safety of Decimus Brutus, and the veterans follow them; for they see that what must be decided by arms is the liberty of the Roman people, not their own advantages.

\ciceroSection{17} What reason is there, then, why the army of Marcus Brutus should be suspect to those who wish Decimus Brutus preserved by every means? Or, indeed, if there were anything that seemed to be feared from Marcus Brutus, would Pansa not see it? or, if he saw it, not be troubled? Who is either wiser in conjecturing things to come, or more diligent in warding off fear? And yet you have seen his disposition toward Marcus Brutus and his zeal. He laid down in his speech what we ought to decree concerning Marcus Brutus and what we ought to think; and so far was he from holding the army of Marcus Brutus dangerous to the commonwealth that in it he placed the firmest and weightiest protection of the commonwealth. No doubt Pansa either does not see this --- for he is of dull wit --- or neglects it: for he does not care that what Caesar did should stand --- he who, on our authority, is about to bring a law before the centuriate assembly to confirm and ratify those acts! Let them cease, then --- either those who feel no fear, from feigning fear and a care for the commonwealth, or those who fear everything, from being too timid --- lest the pretence of the one and the cowardice of the other do us harm. What, confound it ---

\ciceroSection{18} is this the method, always to set the name of the veterans against the best of causes? Even though I embrace their virtue, as I do, yet if they were arrogant I could not endure their disdain. But shall it hinder us, as we try to break the chains of slavery, if someone says the veterans are unwilling? For there are not, I suppose, countless men who would take up arms for the common liberty; there is no man, except the veteran soldiers, who is roused by a freeborn indignation to drive off slavery; the commonwealth, then, can stand relying on the veterans without a great reinforcement of the youth! These, indeed --- the helpers of liberty --- you ought to embrace: the authors of slavery you ought not to follow.

\ciceroSection{19} Finally --- for let a true word, and one worthy of me, burst out at last! --- if the minds of this order are to be steered at the nod of the veterans, and all our words and deeds are to be referred to their will, then death is to be wished for, which to Roman citizens has always been preferable to slavery. All slavery is wretched; but grant that some slavery was unavoidable: what beginning of laying hold of liberty could be more to be desired? Or, when we did not endure that necessary and almost fated misfortune, shall we endure this voluntary one? All Italy has blazed up with longing for liberty; the state can be a slave no longer; later did we give the Roman people this war-garb and these arms than they were demanded of us by the people.

\ciceroSection{20} With great hope, indeed, and one almost assured, we have taken up the cause of liberty; but, granting that the outcomes of war are uncertain and that Mars belongs to both sides alike, yet for liberty we must fight it out at the hazard of life. For life does not consist in breath; rather, there is none at all for one who is a slave. All other nations can bear slavery: our state cannot, and for no other reason than this --- that they flee toil and pain, and to be rid of these can endure all things, whereas we have been so trained and steeped by our ancestors that we refer all our counsels and deeds to dignity and to virtue. So glorious is the recovery of liberty that not even death is to be shunned in the reclaiming of liberty. And if immortality were to follow upon escape from the present danger, even so that escape would seem the more to be shunned, the longer the slavery it bought. But since, day and night, the fates on every side beset us, it is not the part of a man, and least of all a Roman, to hesitate to render back to his country the breath he owes to nature.

\ciceroSection{21} From every side men rush together to put out the common conflagration; the veterans first, following the authority of Caesar, drove back the attempt of Antony; then the Martian legion broke his fury, the Fourth crushed it. Thus, condemned by his own legions, he burst into Gaul, which he found hostile and unfriendly to him in arms and in spirit. The armies of Aulus Hirtius and Gaius Caesar pursued him; afterwards the levy of Pansa roused the city and all Italy; he alone of all men is the enemy. Though he has with him, to be sure, his brother Lucius --- a citizen most dear to the Roman people, whose absence the state can no longer bear.

\ciceroSection{22} What is more foul than that beast, what more monstrous? --- a man who seems to have been born for this very reason, that Marcus Antonius might not be the basest of all mortals. There is, besides, Trebellius, who now, with a cancellation of debts, comes back into favour; Titus Plancus and the rest, his peers: men who fight for this, who labour for this, to be seen restored in defiance of the commonwealth. And the Saxas and the Cafos stir up ignorant men --- themselves boors and rustics, who have never seen this commonwealth and do not wish to see it set in order, who defend not the acts of Caesar but of Antony; whom the boundless possession of the Campanian land has led astray --- and I marvel that they feel no shame at that possession, when they see that they have actors and actresses for their neighbours.

\ciceroSection{23} To crush these plagues, why should we take it amiss that the army of Marcus Brutus has been added? Of an immoderate and turbulent man, I suppose: take care, rather, lest he be one almost too forbearing. And yet in that man's counsels and deeds there was never anything either too much or too little. The whole will of Marcus Brutus, senators, his whole thought, his entire mind, looks to the authority of the Senate and the liberty of the Roman people: these he holds before him, these he wishes to guard. He tried what forbearance could accomplish: when it availed nothing, he judged that force must be met with force. And to him, senators, you ought now to grant the same that, on the thirteenth day before the Kalends of January, on my motion, you granted to Decimus Brutus and Gaius Caesar --- men whose private resolve and action on behalf of the commonwealth was approved and praised by your authority.

\ciceroSection{24} The same you ought to do in the case of Marcus Brutus, by whom an unhoped-for and sudden protection for the commonwealth --- great and firm forces of legions, of cavalry, of auxiliaries --- has been provided. To him must be joined Quintus Hortensius, who, while he held Macedonia, showed himself a most faithful and most steadfast helper to Brutus in raising the army. For concerning Marcus Apuleius I hold that a separate motion should be brought, to whom Marcus Brutus bears witness by letter that he was foremost in the effort of raising the army. Since these things are so, on the matter that Gaius

\ciceroSection{25} Pansa the consul has raised concerning the letter that was brought from Quintus Caepio Brutus, proconsul, and read out before this order, on that matter I move thus: ``Whereas, by the effort, the counsel, the industry, and the virtue of Quintus Caepio Brutus, proconsul, at a most difficult time for the commonwealth, the province of Macedonia and Illyricum and all Greece, and the legions, the army, the cavalry, are in the power of the consuls, the Senate, and the Roman people, Quintus Caepio Brutus, proconsul, has acted well and for the good of the commonwealth, in keeping with his own dignity and that of his ancestors and with the tradition of governing the commonwealth well; and this is, and shall be, welcome to the Senate and the Roman people; and that Quintus

\ciceroSection{26} Caepio Brutus, proconsul, shall guard, defend, watch over, and keep safe the province of Macedonia, Illyricum, and all Greece, and shall command the army which he himself has established and raised; and shall levy and use, for military purposes, if any be needed, such money as is public and can be exacted; and shall take loans of money for military purposes from whomever shall seem good; and shall requisition grain; and shall see to it that, with his forces, he be as near to Italy as possible; and whereas it has been understood from the letter of Quintus Caepio Brutus, proconsul, that the commonwealth has been vigorously aided by the effort and virtue of Quintus Hortensius, proconsul, and that all his counsels were joined with the counsels of Quintus Caepio Brutus, proconsul, and that this has been of great service to the commonwealth, Quintus Hortensius, proconsul, has acted rightly and in due order and for the good of the commonwealth; and it is the pleasure of the Senate that Quintus Hortensius, proconsul, with his quaestor or proquaestor and his own legates, shall hold the province of Macedonia until a successor be appointed to him by decree of the Senate.''
